\subsection{A choice of large repaired targets with persistent second homology}

\begin{lemma}[Proved: convexity after a sufficiently long attachment]
\label{C:large-target-convexity}
Let $f,g:S\rightrightarrows T$ be a parallel pair at a finite stage of
the repaired AKR construction. Suppose that each of the induced graph
embeddings $G_S\to G_T$ is injective on connected components and that
the image of each component is geodesically convex. Here geodesically
convex means that every shortest graph path between two of its
vertices lies in that image. Let $D$ be the largest diameter of a
component of the finite graph $G_T$, taking the diameter of a singleton
to be zero. Choose the repaired path bound $n\geq2$ and a parameter
$s=4k+1$, $k\geq n$, such that
\[
 s>(2n-2)(D+1).
 \tag{LC}\label{eq:akr-large-convexity-bound}
\]
Then every generating map $\alpha_i:T\to Z$ of the repaired target
has geodesically convex image on each component of $G_T$.
This lemma alone does not assert injectivity on components of
$G_T\to G_Z$.
\end{lemma}

\begin{proof}
We first record the elementary linear-amalgam fact used below.
If two graphs are glued along a graph whose maps into both pieces
are injective on components and have geodesically convex component
images, then each component of either piece is geodesically convex
in the amalgam, and the two piece maps are injective on components.
Indeed, a component of either piece contains at most one component
of the gluing graph. Thus the component incidence graph of this
gluing is a matching: each nontrivial resulting component is one
component of each piece glued along a single connected convex
subgraph. Any excursion of a shortest path from one piece into the
other has endpoints in that common subgraph. If the excursion were
shortest between its endpoints, convexity in the other piece would
place it in the common subgraph; if it were not shortest, replacing
it would shorten the original path. Either possibility rules out a
shortest excursion outside the original piece. The component and
convexity claims follow. Induction applies this fact to a linear
sequence of copies of $G_T$, glued along copies of $G_S$ by any of
the neighboring pairs $(g,f),(g,g),(f,g),(f,f)$. Every copy is
componentwise convex and injective on components in that linear
amalgam. At each induction step the image of the next gluing
component is convex in the previous amalgam by transitivity of
convexity, and is still injective on components.

Now fix a generating copy, two vertices in one of its source
components, and a shortest path $P$ between their images in $G_Z$.
The path in the old component gives $|P|\leq D$. As in the cell-fiber
proof of Proposition~\ref{C:fixed-path-kernel}, a target vertex is
visible either in every copy, or in a proper cyclic interval of
indices. The explicit repaired path formulas bound the size of
every proper such interval by $2n-2$; a private point has a singleton
interval. Explicitly, on an actual path of $l\geq2$ vertices the
two middle visibility ranges have size $l$, and the two folded
ranges have sizes $2(u+1)$ for $0\leq u\leq l-2$ and
$2(4k-u+1)$ for $4k-l+2\leq u\leq4k$. Each is at most
$2l-2\leq2n-2$. There are at most $D+1$ vertices of $P$. By
\eqref{eq:akr-large-convexity-bound}, the union of all proper
visibility intervals of these vertices misses an index $r$.
Vertices visible in every copy are deliberately omitted from this
union.

Delete the one gluing overlap from copy $r-1$ to copy $r$, with indices
taken modulo $s$. The remaining diagram is a linear chain containing
all $s$ copies of $G_T$. Let $L$ be its graph colimit. Its canonical
map $q:L\to G_Z$ is the quotient which restores the omitted overlap.
For every vertex of $P$, its cell-fiber graph remains connected after
that overlap is deleted: a proper visibility interval avoids $r$,
so it has no edge crossing this cut; a full circular fiber becomes
a path. The same is true for every edge of $P$. Its visibility set
is an interval or the entire index cycle. If either endpoint has
proper visibility, the edge's interval lies in that endpoint's
interval and avoids $r$. If both endpoints have full visibility,
the routing invariant $(c_{m+1})$ factors the edge through every
copy, since its two endpoints lie in that copy's image. Its fiber
is circular and again becomes a path.

The preimages of a cell of $G_Z$ in $L$ are exactly the components of
that cell's fiber after the indicated gluing edge is removed.
Thus each vertex and edge used by $P$ has a unique preimage in $L$.
These preimages respect incidence, because the quotient is a graph
map and the endpoint preimages are unique. They give a path
$\widetilde P$ of the same length in $L$, whose endpoints are the
original vertices of the selected copy of $G_T$.

The selected copy component is convex in $L$ by the linear-amalgam
fact. Its intrinsic distance between those endpoints is therefore
their distance in $L$. This distance is at most $|\widetilde P|$.
On the other hand, mapping an intrinsic shortest path to $G_Z$
shows $|P|$ is at most that same intrinsic distance. The two lengths
are equal, so $\widetilde P$ is shortest in $L$ and convexity places
it entirely in the selected copy. Mapping back by $q$ places $P$
in that copy's image. This proves the claim for every shortest path.
\end{proof}

\begin{theorem}[Proved: a selected repaired AKR category is not total]
\label{C:selected-nontotal}
For every sequence $(a_m)$ satisfying the numerical requirements of
the repaired AKR Construction C.1, one can choose its finite targets
so that all of the following hold at every stage:
\begin{enumerate}
\item the underlying point sets $X$, and separately the sets
      $H(X)=\pi_0G_X$, admit linear orders for which every arrow
      induces an increasing injection on each of these two sets;
\item every arrow image component in $G_X$ is geodesically convex;
\item every attachment induces an injection on the second integral
      homology of the binary category of elements.
\end{enumerate}
For the resulting countable union $A=\bigcup_mK_m$, the action
\[
 H=h_V/(bv_1\sim bv_2)
\]
satisfies
\[
 H_2(B\operatorname{El}_A H;\mathbb Z)\ne0.
\]
In particular there exists a quasi-filtered category which is not
total. This is an existence theorem about selected sufficiently
large targets; it does not assert the universal choice statement
of Part C.1.
\end{theorem}

\begin{proof}
Initially $G_V$ is a point and $G_W$ is one edge and one isolated
vertex. Give their underlying point sets and their component sets
arbitrary linear orders. Every initial arrow is then increasing
on points and on components, injective on components, and
componentwise convex. The point order and component order are
auxiliary independent orders; the quotient map $X\to H(X)$ is not
required to preserve them.

Assume these assertions at a finite stage $K_m$, and consider one
parallel pair $f,g:S\rightrightarrows T$. By the induction hypothesis
both the point maps $f,g:S\rightrightarrows T$ and the component maps
$H(f),H(g):H(S)\rightrightarrows H(T)$ are increasing injections of
finite linear orders. Neither partial transition relation has a
nontrivial directed cycle: on the finite set of vertices of such a
cycle the induced increasing bijection would be an increasing
permutation, hence the identity. Both relations therefore consist
of fixed points and finite directed paths. More generally every
closed morphism-chain relation is a composition of restrictions
of increasing maps and their increasing partial inverses. It is
an increasing partial bijection on a finite order and has the
same boundedness property. Thus the AKR property (4) holds here
directly from these auxiliary orders.

Choose an odd integer $n\geq3$ at least as large as the number of
vertices in every nontrivial point path and every nontrivial
component path. Apply Proposition~\ref{C:ordered-necklace} to the ordered
component pair and to the underlying ordered point pair. For each
it permits every sufficiently large $k$ and orders the strict set
colimit of its repaired necklace so that all canonical copies are
increasing injections. Choose $k$ beyond both thresholds, and still
larger if necessary so that
\eqref{eq:akr-large-convexity-bound} holds and so that
\[
 s=4k+1,\qquad s/(3n)\in\mathbb N,\qquad
 \frac{s}{3n}\geq\frac{1}{1-a_{m+1}/a_m}.
\]
This is possible: since $3n$ is odd, the integers $s$ divisible by
$3n$ and congruent to $1$ modulo $4$ form an unbounded arithmetic
progression. The ordinary finite repaired target with these
parameters is an admissible target of C.1. Its point-set order is
the order supplied by Proposition~\ref{C:ordered-necklace} and
Lemma~\ref{C:necklace-set-colimit} for the point pair; all $\alpha_i$ are
increasing as actual point maps.

We check that the component order just chosen is an order on the
actual $H(Z)$. The vertex and edge fiber calculation in
Proposition~\ref{C:fixed-path-kernel} identifies $G_Z$ with the
strict graph colimit of the neighboring diagram of $G_S$ and
$G_T$. Graph components are computed by the coequalizer of the two
endpoint maps from the edge set to the vertex set. Since colimits
of sets commute with colimits, taking components preserves this
graph colimit. Hence $H(Z)$ is the strict set colimit of the same necklace
of $H(S),H(T)$. Lemma~\ref{C:necklace-set-colimit}, applied to their
actual component path sizes, identifies this with the component
target ordered above. Transfer that order along this canonical
bijection. The maps $H(\alpha_i)$ are increasing injections.

Lemma~\ref{C:large-target-convexity} supplies convexity of every
component image of $\alpha_i$. Every other incoming arrow is
$\alpha_i b$ for an old arrow $b$, so increasing component
injectivity follows by composition, actual point monotonicity
also follows by composition, and componentwise convexity follows
by transitivity. For a weak disjoint sum use the disjoint union
graph and concatenate the two old orders, separately for points
and for components. Its
incoming maps satisfy both assertions. There are only finitely
many pairs and sums at a finite stage; make these choices for all
of them, using fresh targets as the construction prescribes.
No maps between distinct new targets are added. This proves the
two geometric assertions at every stage.

\paragraph{A deterministic instance.}
To remove all choices from one counterexample, fix
\[
 a_0=2,\qquad
 a_m=\frac32+\frac{1}{10\cdot2^{m-1}}\quad(m\geq1).
\]
These numbers strictly decrease and remain above $3/2$. The required
reciprocals are $5$ at $m=0$ and $15\cdot2^m+2$ at $m\geq1$, hence
integers. Initially label $V$ and $W$ by $0,1$, label their points by
integers, order the three points of $W$ by $v_1<v_2<v_3$, and order its
two graph components with the edge component first. Enumerate all
subsequent objects, points, arrows, and ordered pairs lexicographically
by their finite construction tags; arrows are stored as their tuples
of point values.

At stage $m$, take the least odd $n\geq3$ bounding every point-path
and component-path size for every parallel pair of the finite stage.
Let $D$ be the largest diameter of any of its graph components.
Test $k=n,n+1,\ldots$ in order. Retain only those satisfying the three
displayed C.1 arithmetic conditions and
$4k+1>(2n-2)(D+1)$. For each retained $k$, form the finite repaired
target for every ordered parallel pair. On each target point set
form the precedence graph with arcs
\[
 \alpha_i(x)\longrightarrow\alpha_i(y)
 \quad\text{for all }i\text{ and all }x<y\text{ in }T;
\]
on its actual component set form the analogous graph with arcs
$H(\alpha_i)(C)\to H(\alpha_i)(C')$ for $C<C'$ in $H(T)$.
A loop is counted as a directed cycle. Choose the first retained
$k$ for which all these finitely many precedence graphs are acyclic.
The two orderability applications prove that this search terminates.
Assign each resulting point set and component set its lexicographically
least topological ordering, always selecting the least available tag
at a tie. This is a prescribed finite algorithm. It can use one common
$n,k$ for all pairs at the stage, so it also covers that more restrictive
parameter convention.

Fixed-point, private-point, and path-coordinate tags specify the
covered repaired target uniquely; path components are indexed by
their least old point tag. Give every target a fresh tag consisting
of its stage, construction type, source and target tags, and ordered
pair of arrow tuples. For each ordered pair of old objects also add
the freshly tagged disjoint-sum target, with left/right point tags
and the concatenated point and component orders. Keep the old
hom-sets, insert the prescribed generating maps, and close under
composition as actual functions, identifying equal function tuples
with the same source and target. These finite rules specify all
objects and arrows. Composition is ordinary function composition,
so it is associative and has the indicated identity functions.
Each stage is finite, hence the union is a countable small category.
The following homology argument applies to this fully specified
instance as well as to the preceding existence construction.

For every parallel pair at such a stage,
Lemma~\ref{C:convex-fixed} now applies: its old images are convex,
its component maps are injective, and actual point monotonicity
excludes a nontrivial point-transition cycle. Its fixed subgraph therefore meets each
source component in an empty or connected subgraph. The exact
kernel formula of Proposition~\ref{C:fixed-path-kernel} shows that
adjoining the pair's quasi-filter target kills no old second
homology class. A disjoint-sum target also kills none: its cycle
group is the direct sum of the cycle groups of its two summands,
so the map $a$ in Proposition~\ref{C:attachment-cycle-test} is
injective. Adding the finitely many new targets one at a time
proves injectivity across every full stage.

Finally the first stage already has a nonzero class, independently
of how large these first choices were. The two quasi-filter targets
$P_+,P_-$ for $(v_1,v_2)$ and $(v_2,v_1)$ each have one binary graph
cycle; orient their generators $\kappa_+,\kappa_-$ so that their
augmentations are $1$, as follows from the signed neighboring count
$2k-(2k-1)=1$. All other first-stage graphs are forests. For a pair
involving $v_3$ and exactly one of $v_1,v_2$, every binary edge has
its other endpoint private to that copy, so it cannot belong to a
cycle. For a repeated point map the binary graph is a star, a
disjoint union of edges, and isolated vertices. The identity
quasi-filters and the initial weak sums also give forests.
There are no arrows between distinct new first-stage targets, and
the old cycle groups at $V,W$ are zero. Consequently
\[
 \operatorname{colim}_{K_1}K=\mathbb Z\kappa_+
                   \oplus\mathbb Z\kappa_-,
 \qquad
 H_2(B\operatorname{El}_{K_1}H;\mathbb Z)
       =\ker\bigl(\mathbb Z^2\xrightarrow{(1,1)}\mathbb Z\bigr)
       \cong\mathbb Z.
\]
The generator $\kappa_+-\kappa_-$ survives every subsequent stage.
Every finite nerve chain of the union occurs in some stage, so its
normalized chain complex is the filtered union of the stage chain
complexes. Filtered unions of abelian groups are exact; therefore
homology commutes with this union, and the persistent generator
is nonzero in the displayed $H_2$ of $A$. The repaired construction
is quasi-filtered, whereas totality would force this category of
elements to have zero positive homology. Hence $A$ has the claimed
properties.
\end{proof}
